By the sharp-frontier theorem, the image of the full response-vector target over \(\mathfrak F_x(P_X)\) is the nonempty set \(\mathcal R_x^{\mathrm{vec}}(C)\). Identification therefore means exactly that this set is a singleton. If two indices \(j,k\in J_x(C)\) produced the same vector, then their nonzero treatment entries would give
\[
 \frac{C_{-x,j}}{C_{xj}}=\frac{C_{-x,k}}{C_{xk}}. \tag{35}
\]
Adjoining the common \(x\)-coordinate \(1\) to both sides of (35) yields \(c_j/C_{xj}=c_k/C_{xk}\), so \(c_j\) and \(c_k\) are proportional. Pairwise nonproportionality forces \(j=k\). Hence the map \(j\mapsto C_{-x,j}/C_{xj}\) is injective on \(J_x(C)\), and
\[
 |\mathcal R_x^{\mathrm{vec}}(C)|=|J_x(C)|. \tag{36}
\]
This proves the vector criterion.

For the scalar criterion, every \(j\in J_x(C)\) has \(C_{xj}\ne0\), so the \(x\)-row of \(C_{\{x,y\},J_x(C)}\) is nonzero. This two-row matrix has rank one exactly when all its two-column determinants vanish:
\[
 C_{xj}C_{yk}-C_{xk}C_{yj}=0
 \quad\text{for all }j,k\in J_x(C). \tag{37}
\]
Division by the nonzero product \(C_{xj}C_{xk}\) makes (37) equivalent to
\[
 \frac{C_{yj}}{C_{xj}}=\frac{C_{yk}}{C_{xk}}
 \quad\text{for all }j,k\in J_x(C), \tag{38}
\]
which is precisely \(|\mathcal R_{xy}(C)|=1\). The scalar part of the sharp theorem identifies this with scalar causal identification.

If \(|J_x(C)|>1\), choose distinct \(j,k\in J_x(C)\). Equation (36) says their response vectors differ. The canonical witness lemma returns \(\mathcal M_j,\mathcal M_k\in\mathfrak F_x(P_X)\) realizing those two vectors. If instead the scalar rank condition fails, (37) supplies a pair whose scalar ratios differ, and the same two witnesses realize the distinct scalar responses. After the shared \(O(n^3)\) basis extension and inverse, each of two shears, complementary matchings, and row normalizations costs \(O(n^3)\), proving the additional-operation claim.